%% kapa: exam=1 round=33 year=2022 period=2 form=A subject=accounting subject_ko=회계학 seq=19 number=59
%% source: Q-Net 감정평가사 기출문제 (cst00309) artlSeq=5209997 file=raw/1008/5209997/2022년도 제33회 감정평가사 1차시험 2교시(감정평가 관계법규, 회계학).pdf
%% conversion: kordoc 4.24.51 -> tools/gen1.py (auto); review=proofread-sonnet; flags=none
\begin{question}{59}
20×1년 초에 사업을 개시한 (주)감평의 회계담당자는 20×1년 말에 정기예금에 대한 미수이자 ￦200을 계상하지 않은 오류를 발견하였다. (주)감평의 당기 및 차기 이후 적용 법인세율이 모두 30\%일 때 이러한 회계처리 오류가 (주)감평의 20×1년도 재무제표에 미친 영향으로 옳은 것은? (단, 세법상 정기예금 이자는 이자수령시점에 과세된다. 또한, (주)감평은 이연법인세 회계를 적용하고 있으며, 미래에 충분한 과세소득이 발생할 것으로 판단된다.)
\begin{choices}
  \item 당기법인세자산이 ￦60 과대계상 되었다.
  \item 당기법인세부채가 ￦60 과소계상 되었다.
  \item 법인세비용이 ￦60 과대계상 되었다.
  \item 당기순이익이 ￦140 과소계상 되었다.
  \item 이연법인세자산이 ￦60 과소계상 되었다.
\end{choices}
\end{question}
\answer{4}
